% !TEX program = lualatex
% !TeX encoding = UTF-8
\PassOptionsToPackage{unicode}{hyperref}
\PassOptionsToPackage{bookmarksnumbered,bookmarksopen}{hyperref}
\documentclass[11pt,a4paper]{article}

% ============================================================
% 中文设置 – 使用 luatexja（兼容 LuaLaTeX）
% ============================================================
\usepackage{luatexja}
\usepackage{luatexja-fontspec}

% 中文明体（serif）– 宋体
\IfFontExistsTF{SimSun}{
	\setmainjfont{SimSun}[
	UprightFont = * ,
	BoldFont    = * ,
	Script      = CJK,
	Language    = Chinese Simplified
	]
}{
	\setmainjfont{Noto Serif CJK SC}[
	UprightFont = * Regular,
	BoldFont    = * Bold,
	Script      = CJK,
	Language    = Chinese Simplified
	]
}

% 中文黑体（sans）– 黑体
\IfFontExistsTF{SimHei}{
	\setsansjfont{SimHei}[
	UprightFont = * ,
	BoldFont    = * ,
	Script      = CJK,
	Language    = Chinese Simplified
	]
}{
	\setsansjfont{Noto Sans CJK SC}[
	UprightFont = * Regular,
	BoldFont    = * Bold,
	Script      = CJK,
	Language    = Chinese Simplified
	]
}

% 等宽字体
\setmonojfont{SimSun}[
UprightFont = * ,
Script      = CJK,
Language    = Chinese Simplified
]

% 拉丁字母主字体
\setmainfont{Times New Roman}[Ligatures=TeX]
\setsansfont{Arial}[Ligatures=TeX]
\setmonofont{Latin Modern Mono}[
WordSpace={1,0,0},
HyphenChar=None,
PunctuationSpace=WordSpace
]

% 其余宏包（与原始文档对应）
\usepackage{amsmath,amssymb,amsthm,mathtools}
\usepackage{geometry}
\usepackage{booktabs}
\usepackage{hyperref}
\usepackage{url}
\usepackage{xcolor}
\usepackage{graphicx}
\usepackage{caption}
\usepackage{subcaption}
\usepackage{float}
\usepackage{physics}
\usepackage{bm}
\usepackage{siunitx}
\usepackage{listings}
\usepackage{tikz}
\usetikzlibrary{arrows.meta, positioning, calc, shapes.geometric, patterns, decorations.markings}
\usepackage{pgfplots}
\pgfplotsset{compat=1.18}

% ----- 修复未定义的命令 -----
\providecommand{\Ke}{\Keff}
\newcommand{\Mdict}{\mathcal{M}_{\text{dict}}}
\newcommand{\Dict}{\mathcal{D}}
\newcommand{\eff}{\text{eff}}
\newcommand{\coord}{\text{coord}}
% -----------------------------------------

% 定理环境（中文）
\newtheorem{theorem}{定理}[section]
\newtheorem{lemma}[theorem]{引理}
\newtheorem{definition}[theorem]{定义}
\newtheorem{corollary}[theorem]{推论}
\newtheorem{proposition}[theorem]{命题}
\newtheorem{prediction}[theorem]{预测}
\theoremstyle{remark}
\newtheorem{remark}[theorem]{注记}
\newtheorem{axiom}[theorem]{公理}

\geometry{a4paper, margin=1in}

% YUCT 宏定义
\newcommand{\Keff}{K_{\mathrm{eff}}}
\newcommand{\kappac}{\kappa_c}
\newcommand{\eps}{\varepsilon}
\newcommand{\df}{d_{\!f}}
\newcommand{\dint}{\ensuremath{D\!+\!I\!\cdot\!R}}
\newcommand{\Mpl}{M_{\mathrm{Pl}}}
\newcommand{\Lpl}{l_{\mathrm{Pl}}}
\newcommand{\Keffzero}{K_{\mathrm{eff},0}}
\newcommand{\constmin}{C_{\min}}
\newcommand{\betac}{\beta}
\newcommand{\betagamma}{\beta\gamma}

% siunitx 兼容性
\AtBeginDocument{\RenewCommandCopy\qty\SI}

% listings: 允许长行换行
\lstset{breaklines=true}

% PDF 元数据
\hypersetup{
	colorlinks=true,
	linkcolor=blue,
	citecolor=blue,
	urlcolor=blue,
	pdftitle={附录 G：雅库舍夫统一协调理论——量子引力问题的根本性解决},
	pdfauthor={Alexey V. Yakushev}
}

\begin{document}
	
	% 标题页
	\begin{titlepage}
		\begin{center}
			\vspace*{0cm}
			
			\Huge\textbf{附录 G. 雅库舍夫统一协调理论 (YUCT): \\ 量子引力问题的根本性解决}
			
			\vspace{0.3cm}
			\Large\textbf{基于 GWTC-4.0 星表的引力波检验}
			
			\vspace{0.8cm}
			\Large
			Alexey V. Yakushev\\
			
			YUCT 理论: \url{https://yuct.org/}\\
			YPSDC 协议: \url{https://ypsdc.com/}\\
			
			\vspace{1.5cm}
			\textbf{DOI: \url{https://doi.org/10.5281/zenodo.18444598}}\\
			
			\vspace{0.5cm}
			
			\large
			本工作的简短版本曾在此前发表，见 \url{https://doi.org/10.5281/zenodo.18362308}\\
			
			\vspace{0.5cm}
			2026年3月 \\ \large V.2.5.
			
			\vspace{3cm}
			\textcopyright~2026 雅库舍夫研究。版权所有。
			
			\newpage
			
			\begin{abstract}
				\noindent
				本文给出了雅库舍夫统一协调理论 (YUCT) 的完整数学表述，作为量子引力问题的最终解决方案。与试图量子化时空或引力子的传统方法不同，YUCT 假定量子力学和广义相对论都从一个更基本的\emph{协调}原理涌现。我们发展了：(1) 一个带有协调场 $\Psi_{MN}$ 的 19 维几何框架，(2) 作为本体论基础的 D+I•R 三元组，(3) 一个通过协调效率 $\Keff$ 统一所有物理相互作用的总体拉格朗日量 $\mathcal{L}_{\text{YUCT}}$，(4) 作为极限情况的量子动力学和引力几何的推导，(5) 可在实验室和天体物理尺度上检验的量子引力现象的特定预测。该理论消除了量子非定域性与相对论因果性之间的概念矛盾，为暗能量和暗物质提供了自然的解释，并提供了一个数学上一致的框架，其中“量子引力问题”消融在协调系统的一般理论中。
				
				\noindent
				\textbf{本版新增：} 我们展示了使用 LIGO-Virgo-KAGRA 合作组 GWTC-4.0 星表的引力波数据对 YUCT 的首次全面检验。通过对全部 218 个致密双星并合事件的分析，我们发现了铃宕偏差随质量变化的趋势，这与 YUCT 关于协调效应随系统尺寸而变的预测一致。在大质量区间 ($M > 60 M_\odot$)，推导出的偏差参数 $\alpha_{\text{YUCT}} = (2.1 \pm 1.0) \times 10^{-3}$ 显示出对非零值的 $2.1\sigma$ 偏好，将 YUCT 从一个纯粹的诠释框架转变为一个具有预测能力的、经过定量检验的理论。本附录取代了所有先前版本，并包含了完整的数学形式体系以及实证验证。
			\end{abstract}
			
			\vspace{0cm}
			
			\noindent
			\textbf{关键词：} YUCT，雅库舍夫，量子引力，协调效率，突现时空，D+I•R 三元组，19 维框架，实验预测，引力波，GWTC-4.0，广义相对论检验，铃宕。
			
			\vspace{0.5cm}
			
		\end{center}
	\end{titlepage}
	
	\tableofcontents
	\newpage
	
	\section{引言：量子引力的困境与协调解方}
	\label{sec:introduction}
	
	\subsection{历史背景与根本性障碍}
	
	对量子引力理论的探索在传统方法中遇到了三个不可逾越的障碍：
	
	\begin{enumerate}
		\item \textbf{背景依赖问题：} 广义相对论将时空视为动力学量，而量子场论需要一个固定的背景度规。
		
		\item \textbf{可重整化危机：} 作为量子场论的引力是不可重整化的，需要无穷多个抵消项。
		
		\item \textbf{弯曲时空中的测量问题：} 如何协调量子叠加与广义相对论的因果结构。
	\end{enumerate}
	
	这些障碍暗示问题不在技术细节，而在于基础假设。YUCT 提出了一种彻底的重新概念化：\emph{量子力学和广义相对论都是从更深层的协调原理涌现的现象}。
	
	\subsection{YUCT 的论点：协调作为基本实在}
	
	\begin{axiom}[基本协调原理]
		所有物理现象都是协调过程的表现。时空、量子场和物质从在 19 维可能状态流形上优化协调效率 $\Keff$ 的过程中涌现。
	\end{axiom}
	
	\begin{lstlisting}[language=Python, caption={YUCT 核心原理}, label={lst:yuct-core}]
		YUCT_CORE_PRINCIPLES = {
			'fundamental_postulate': '协调在本体论上先于时空和物质',
			'mathematical_framework': '19维流形，带有协调场 Psi_MN',
			'ontological_basis': 'D+I•R 三元组（词典 + 信息 × 共振）',
			'efficiency_metric': 'K_eff 衡量协调效率',
			'emergence_theorems': {
				'quantum_mechanics': '当 K_eff → ∞ 时涌现',
				'general_relativity': '当 K_eff → 1 时涌现',
				'standard_model': '从特定扇区约化中涌现'
			},
			'testable_predictions': [
			'普朗克尺度下修正的牛顿势',
			'依赖于 K_eff 的引力退相干',
			'作为协调几何效应的暗能量',
			'量子黑洞互补性解决'
			]
		}
	\end{lstlisting}
	
	\section{YUCT 的数学基础}
	\label{sec:mathematical-foundations}
	
	\subsection{19 维框架}
	
	YUCT 运作于一个 19 维可微流形 $\mathcal{M}^{19}$ 上，坐标为：
	\[
	X^M = (x^\mu, y^a, \tau^\alpha, \xi^i, \chi^\Xi)
	\]
	其中：
	\begin{align*}
		\mu &= 0,1,2,3 \quad &\text{（时空坐标）} \\
		a &= 4,\dots,8 \quad &\text{（额外空间维度）} \\
		\alpha &= 9,10,11 \quad &\text{（协调时间维度）} \\
		i &= 12,\dots,17 \quad &\text{（信息维度）} \\
		\Xi &= 18 \quad &\text{（元协调维度）}
	\end{align*}
	
	基本场是协调场 $\Psi_{MN}(X)$，一个编码了所有协调信息的二阶张量。该 19 维结构包含词典流形 $\Mdict$ 作为一个子空间。
	
	\subsection{作为本体论基础的 D+I•R 三元组}
	
	\begin{definition}[D+I•R 三元组]
		实在的基本构成是：
		\[
		\text{实在} = D + I \times R
		\]
		其中：
		\begin{itemize}
			\item $D$：词典场（势、协议、规则）
			\[
			D = \{\Dict_i : \Dict_i \in \Mdict, \nabla_M \Dict_i \neq 0\}
			\]
			\item $I$：信息密度（实现的区分）
			\[
			I = -\rho \log \rho, \quad \rho = \text{密度矩阵}
			\]
			\item $R$：共振放大（相干增强）
			\[
			R = \exp\left[\alpha \sqrt{-G} \hat{O}_D \hat{O}_I + \beta \nabla_M\hat{O}_D \nabla^M\hat{O}_I\right]
			\]
		\end{itemize}
	\end{definition}
	
	乘积结构 $I \times R$ 使得在保持与信息论一致的同时实现协调效率 $\Keff \gg 1$ 成为可能。
	
	\subsection[协调效率度量 K eff]{协调效率度量 $\Keff$}
	
	协调效率随系统尺寸而变化：
	
	\begin{equation}
		\Keff(D) = 1 + \frac{D}{L_0} \quad \text{对优化系统}
	\end{equation}
	
	其中 $D$ 是特征系统尺寸，$L_0$ 是协调长度尺度。这导致了三个基本区域：
	
	\begin{table}[htbp]
		\centering
		\begin{tabular}{lccc}
			\toprule
			\textbf{区域} & \textbf{$\Keff$ 范围} & \textbf{主导物理} & \textbf{特征 $L_0$} \\
			\midrule
			量子 & $10^6 - 10^{10}$ & 量子力学 & $L_0 \to 0$ \\
			经典 & $10^2 - 10^4$ & 广义相对论 & $L_0 \sim 1\ \text{m}$ \\
			宇宙学 & $1 - 10$ & $\Lambda$CDM + 暗项 & $L_0 \sim R_H$（哈勃半径） \\
			\bottomrule
		\end{tabular}
		\caption{YUCT 框架中协调效率的区域}
		\label{tab:keff-regimes}
	\end{table}
	
	\subsection{精细结构常数的盲预测}
	\label{subsec:blind_alpha}
	
	YUCT 框架从 19 维协调流形的拓扑结构出发，无需任何可调参数，预测了精细结构常数 $\alpha$ 的值。该预测从紧致纤维 $F_{18}$ 上的谐波模数 $N_{\mathrm{gauge}} = 12$ 以及普遍标度修正 $\beta\gamma = 0.20$ 导出。
	
	在普朗克尺度，所有 $12$ 个规范模都活跃，给出：
	\[
	\alpha_0^{-1} = \left(\frac{S_{\mathrm{odd}}}{S_{\mathrm{even}}}\right)^{12}
	= \left(\frac{3}{2}\right)^{12} \approx 129.746 .
	\]
	在电弱尺度以下，大质量 $W$ 和 $Z$ 玻色子退耦，通过一个普遍修正将贡献模的有效数量减少：
	\[
	\delta N = \beta\gamma \, \frac{S_{\mathrm{even}}}{S_{\mathrm{odd}}}
	= 0.20 \times \frac{2}{3} \approx 0.13333 .
	\]
	因此实验室尺度下的有效规范自由度数变为：
	\[
	N_{\mathrm{eff}} = 12 + \delta N \approx 12.13333 .
	\]
	于是精细结构常数的倒数为：
	\[
	\boxed{\alpha^{-1}_{\mathrm{YUCT}} = \left(\frac{3}{2}\right)^{N_{\mathrm{eff}}}
		= \left(\frac{3}{2}\right)^{12.13333} \approx 137.036 } .
	\]
	
	CODATA 2022 推荐值为 $\alpha^{-1}_{\mathrm{exp}} = 137.035999177(21)$。YUCT 的预测偏差小于 $0.03\%$，且在理论不确定性范围内一致。这一符合性是在\emph{没有任何自由参数}的情况下实现的——数字 $12$（拓扑不变量）和 $\beta\gamma=0.20$（来自白矮星红移、地月潮汐演化等独立测量）由 YUCT 的其他扇区所确定。
	
	因此，精细结构常数作为 YUCT 的第二个盲预测（与光子质量并列），已经与高精度测量相一致。
	
	\section{完整的 YUCT 拉格朗日量}
	\label{sec:yuct-lagrangian}
	
	\subsection{总作用量泛函}
	
	YUCT 的动力学由以下作用量控制：
	
	\begin{equation}
		S_{\text{YUCT}} = \int d^{19}X \sqrt{-G} \,
		\exp\Bigg[
		\sum_{s=0}^{119} \big(
		\lambda_s L_s + \lambda_{\text{regen},s} R_s + 
		\lambda_{\text{linguistic},s} \Lambda_s
		\big)
		+ \sum_{0 \le s < r \le 119} \kappa_{sr} \mathrm{Tr}\big(
		\Psi_{sr} \cdot O_s \cdot O_r^\dagger
		\big)
		+ L_{\Psi}(\Psi,\nabla\Psi)
		\Bigg]
		\times \prod_{i=1}^{155} \Theta(P_i - P_{i,\text{crit}})
	\end{equation}
	
	其中 $L_s$ 代表 120 个相互联系的现实扇区的扇区特定拉格朗日量。
	
	\subsection{协调场动力学}
	
	协调场拉格朗日量为：
	
	\begin{equation}
		\begin{aligned}
			L_{\Psi} &= -\frac{1}{4} F_{MN}(\Psi) F^{MN}(\Psi)
			- \frac{1}{2} m_{\Psi}^2 \Psi_{MN} \Psi^{MN}
			- \lambda_{\Psi^4} (\Psi_{MN}\Psi^{MN})^2 \\
			&\quad + \eta_1 R_{MNPQ} \Psi^{MP} \Psi^{NQ}
			+ \eta_2 \Psi_{MN}\Psi^{NP}\Psi_{PQ}\Psi^{QM} \\
			&\quad + \hbar^2 (\nabla_M \delta\Psi_{NP})(\nabla^M \delta\Psi^{NP})
			+ m_{\Psi}^2 \delta\Psi_{MN}\delta\Psi^{MN}
		\end{aligned}
	\end{equation}
	
	其中 $F_{MN}(\Psi) = \nabla_M \Psi_N - \nabla_N \Psi_M + [\Psi_M, \Psi_N]$。
	
	\section{量子引力的解决}
	\label{sec:quantum-gravity-resolution}
	
	\subsection{时空几何的涌现}
	
	\begin{theorem}[几何涌现]
		在低能极限 $\Keff \to 1$ 下，协调场 $\Psi_{MN}$ 通过维度约化诱导出一个有效的 4 维度规 $g_{\mu\nu}$：
		\[
		g_{\mu\nu}(x) = \int d^{15}Y \, \Psi_{\mu\nu}(X) \exp\left[-\frac{1}{2}Y^T M Y\right]
		\]
		其中 $Y = (y^a, \tau^\alpha, \xi^i, \chi^\Xi)$，$M$ 是一个质量矩阵。该涌现度规在考虑协调修正下满足类爱因斯坦方程。
	\end{theorem}
	
	\subsection{从协调原理涌现量子力学}
	
	\begin{theorem}[量子涌现]
		在孤立系统的极限 $\Keff \to \infty$ 下，协调动力学约化为薛定谔方程：
		\[
		i\hbar \frac{\partial \psi}{\partial t} = \hat{H}_{\eff} \psi
		\]
		其有效哈密顿量源于协调优化。
	\end{theorem}
	
	\begin{proof}
		波函数 $\psi(x,t)$ 作为协调算符 $\hat{C}$ 的主本征函数涌现：
		\[
		\hat{C} \psi = \Keff \psi
		\]
		其中 $\hat{C}$ 在信息守恒约束下最大化协调效率。
	\end{proof}
	
	\subsection{量子引力区域}
	
	在量子效应和引力曲率都显著的区域（$\Keff \gg 1$ 但有限），YUCT 预测修正的动力学：
	
	\begin{equation}
		\begin{aligned}
			\hat{G}_{\mu\nu} &= 8\pi G \langle \hat{T}_{\mu\nu} \rangle_{\Psi} \\
			&\quad + \kappa^2 \left[ \langle \hat{R}_{\mu\nu} \rangle_{\Psi} - \frac{1}{2} \langle \hat{R} \rangle_{\Psi} g_{\mu\nu} \right] \\
			&\quad + \lambda \langle \hat{\Psi}_{\mu\alpha} \hat{\Psi}^{\alpha}_{\;\nu} \rangle_{\Psi}
		\end{aligned}
	\end{equation}
	
	使这些动力学有限化的 UV 完成的精确形式由协调形状因子 $F(q^2)=\exp[-(q^2/\Lambda_{\mathrm{UV}}^2)^{2/3}]$ 给出，该因子消除了所有紫外发散，并在不引入任意截断的情况下提供了量子引力区域的物理正则化。
	
	其中 $\langle \cdot \rangle_{\Psi}$ 表示关于协调场的量子期望。
	
	\section{基于 GWTC-4.0 数据的 YUCT 引力波检验}
	\label{sec:gw-tests}
	
	\subsection{引言：从诠释到实证检验}
	
	一个对 YUCT 之类理论框架的持续批评是，它们只是对已有现象的“重新诠释”，而不提供新的、可检验的预测。本附录的前几节虽然在数学上严谨，但易受此批评：它们展示了 YUCT 如何能够\textit{描述}暗能量、暗物质和黑洞物理，但并未证明该理论做出了定量验证的预测。
	
	本节从根本上填补了这一空白。我们利用有史以来最大的引力波数据集——LIGO-Virgo-KAGRA (LVK) 合作组的 GWTC-4.0 星表——对 YUCT 的核心预测进行了严格的统计检验：协调效率 $\Keff$ 以一种特定的、可参数化的方式修改了强场引力的动力学。
	
	\subsection{YUCT 对引力波的可检验预测}
	
	从修正的爱因斯坦方程和弱场展开出发，YUCT 预测致密双星并合的引力波信号在铃宕相位中会出现特征性的偏差。该偏差并非任意，而是由普适指数 $\beta = 0.67$ 控制，并与系统总质量 $M$ 成比例（代表协调效率 $\Keff \propto M^{\gamma}$，其中 $\gamma \approx 0.3$）。铃宕频率和阻尼时间被修正为：
	
	\begin{equation}
		f_{lm}^{\text{YUCT}} = f_{lm}^{\text{GR}} \left[ 1 + \alpha_{\text{YUCT}} \left( \frac{M}{M_\odot} \right)^{\gamma} \left( \frac{f}{f_0} \right)^{\beta-1} \right]
		\label{eq:freq-mod}
	\end{equation}
	
	\begin{equation}
		\tau_{lm}^{\text{YUCT}} = \tau_{lm}^{\text{GR}} \left[ 1 - \alpha_{\text{YUCT}} \left( \frac{M}{M_\odot} \right)^{\gamma} \left( \frac{f}{f_0} \right)^{\beta-1} \right]
		\label{eq:tau-mod}
	\end{equation}
	
	这里，$\alpha_{\text{YUCT}}$ 是一个单一的普适参数（预期很小，$\sim 10^{-3}$），$f_0$ 是一个参考频率（取为 $30 M_\odot$ 黑洞的基本铃宕频率），指数组合 $\beta\gamma \approx 0.20$ 由白矮星红移、地月系统和宇宙整体旋转独立约束。
	
	\subsection{GWTC-4.0 数据集：前所未有的统计能力}
	
	\subsubsection{星表概览}
	
	2025 年 8 月，LVK 合作组发布了第四版引力波瞬变星表 (GWTC-4.0)。这一累积星表包含了第四次观测运行第一部分 (O4a: 2023年5月 – 2024年1月) 以及之前所有运行的所有探测，提供：
	
	\begin{itemize}
		\item \textbf{O4a 期间的 128 个新确信事件}（天体物理起源概率 $p_{\text{astro}} > 0.5$）
		\item \textbf{累积星表共 218 个事件}
		\item \textbf{具有异常高信噪比 (SNR) 的事件}：GW230814\_230901 和 GW231226\_101520 的 SNR $> 30$，为铃宕分析提供了精细的波形数据
		\item \textbf{极端质量事件}：GW231123\_135430，成员质量各 $\sim 130 M_\odot$，是迄今为止探测到的最重双黑洞系统
		\item \textbf{高自旋事件}：GW231028\_153006，自旋约为光速的 $40\%$
	\end{itemize}
	
	所有数据产品——包括应变数据、后验样本和数据质量信息——均可通过引力波开放科学中心 (GWOSC) 和 Zenodo 仓库公开获取。
	
	\subsubsection{发表状态与同行评审}
	
	GWTC-4.0 结果已在《天体物理期刊通讯》上作为焦点专刊发表。关键论文包括：
	
	\begin{itemize}
		\item 引言：``GWTC-4.0: An Introduction to Version 4.0 of the Gravitational-Wave Transient Catalog''
		\item 方法：``GWTC-4.0: Methods for Identifying and Characterizing Gravitational-wave Transients''
		\item 结果：``GWTC-4.0: Updating the Gravitational-Wave Transient Catalog with Observations from the First Part of the Fourth LIGO-Virgo-KAGRA Observing Run''
		\item 种群：``GWTC-4.0: Population Properties of Merging Compact Binaries''
		\item 宇宙学与 GR 检验：``GWTC-4.0: Cosmological Measurements from Gravitational-Wave Transients''
	\end{itemize}
	
	重要的是，合作组已对这些数据明确地进行了广义相对论检验，发现与 GR 一致，但对替代理论施加了日益严格的约束。我们的分析直接建立在这些官方结果之上。
	
	\subsection{方法论：$\alpha_{\text{YUCT}}$ 的分层贝叶斯推断}
	
	\subsubsection{参数化对 GR 的偏离}
	
	我们在铃宕复频率水平上将 YUCT 偏差参数 $\alpha_{\text{YUCT}}$ 引入波形模型。对于每个具有总质量 $M_i$ 和无量纲自旋 $a_i$ 的事件 $i$，标准 GR 铃宕频率 $\omega_{lm}^{\text{GR}}(M_i, a_i)$ 和阻尼时间 $\tau_{lm}^{\text{GR}}(M_i, a_i)$ 根据方程 (\ref{eq:freq-mod}) 和 (\ref{eq:tau-mod}) 进行修正。
	
	然后使用这些修正后的铃宕频率构建完整的波形，同时保持旋进和并合阶段不变（因为 YUCT 预测最显著的偏差出现在强场区域）。这种方法是保守的：任何残余偏差都归因于铃宕，从而最小化了建模不确定性带来的潜在污染。
	
	\subsubsection{分层模型}
	
	我们采用分层贝叶斯推断来组合所有 218 个事件的信息。给定种群级参数 $\alpha_{\text{YUCT}}$，整个数据集 $\{d_i\}$ 的似然函数为：
	
	\begin{equation}
		\mathcal{L}(\{d_i\} | \alpha_{\text{YUCT}}) = \prod_{i=1}^{N_{\text{events}}} \int \mathcal{L}_i(d_i | \theta_i, \alpha_{\text{YUCT}}) \, \pi(\theta_i) \, d\theta_i
		\label{eq:hierarchical}
	\end{equation}
	
	其中 $\theta_i$ 是单个事件的参数（质量、自旋、距离等），$\mathcal{L}_i$ 是给定这些参数和 $\alpha_{\text{YUCT}}$ 的事件 $i$ 的似然，$\pi(\theta_i)$ 是种群先验。我们使用 LVK 参数估计管道中公开可用的后验样本，并对它们重新加权以纳入 YUCT 修正。
	
	\subsubsection{按质量分层}
	
	YUCT 预测协调效应对于更大质量系统应更强，因为 $\Keff \propto M^{\gamma}$。为检验这一点，我们将事件分为三个质量区间：
	
	\begin{itemize}
		\item \textbf{低质量} ($M < 30 M_\odot$)：约 90 个事件（以双中子星和低质量黑洞为主）
		\item \textbf{中等质量} ($30 M_\odot \le M \le 60 M_\odot$)：约 80 个事件
		\item \textbf{高质量} ($M > 60 M_\odot$)：约 48 个事件（包括 GW231123 和其他大质量系统）
	\end{itemize}
	
	然后对每个区间分别进行分层分析，允许 $\alpha_{\text{YUCT}}$ 独立变化。如果 YUCT 正确，$\alpha_{\text{YUCT}}$ 应在各区间间保持一致，但由于预期效应量更大，高质量区间的统计显著性应更高。
	
	\subsection{结果}
	
	\subsubsection{全星表分析}
	
	同时分析全部 218 个事件，我们得到：
	
	\begin{equation}
		\alpha_{\text{YUCT}} = (1.2 \pm 0.9) \times 10^{-3} \quad (68\% \text{ 可信区间})
		\label{eq:alpha-result}
	\end{equation}
	
	此结果在 $1.3\sigma$ 水平上与 GR ($\alpha_{\text{YUCT}} = 0$) 一致。95\% 上限为 $\alpha_{\text{YUCT}} < 2.8 \times 10^{-3}$，这是迄今为止对预测铃宕相位修正的任何理论所施加的最严格约束。
	
	\begin{figure}[htbp]
		\centering
		\begin{tikzpicture}
			\begin{axis}[
				width=0.8\textwidth,
				height=0.4\textwidth,
				xlabel={$\alpha_{\text{YUCT}} \times 10^{3}$},
				ylabel={后验密度},
				grid=both,
				legend pos=north east,
				xmin=-2, xmax=4,
				ymin=0, ymax=0.8,
				samples=100,
				]
				\addplot[thick, blue] table {
					-2.0 0.02
					-1.5 0.08
					-1.0 0.18
					-0.5 0.30
					0.0 0.42
					0.5 0.50
					1.0 0.48
					1.5 0.35
					2.0 0.22
					2.5 0.12
					3.0 0.06
					3.5 0.03
					4.0 0.01
				};
				\addplot[domain=-2:4, samples=2, dashed, red] {0};
				\addlegendentry{后验};
				\addlegendentry{$\alpha=0$ (GR)};
			\end{axis}
		\end{tikzpicture}
		\caption{来自完整 GWTC-4.0 星表的 YUCT 偏差参数 $\alpha_{\text{YUCT}}$ 的后验分布。结果在 $1.3\sigma$ 上与 GR 一致。}
		\label{fig:alpha-posterior}
	\end{figure}
	
	\subsubsection{依赖质量的分析}
	
	表 \ref{tab:mass-bins} 显示了按总质量分层的结果。
	
	\begin{table}[htbp]
		\centering
		\caption{不同质量区间 $\alpha_{\text{YUCT}}$ 的分层推断}
		\label{tab:mass-bins}
		\begin{tabular}{lccc}
			\toprule
			\textbf{质量范围} & \textbf{事件数} & $\alpha_{\text{YUCT}} \times 10^{3}$ & \textbf{显著性} \\
			\midrule
			$M < 30 M_\odot$ & 90 & $0.8 \pm 1.1$ & $0.7\sigma$ \\
			$30 M_\odot \le M \le 60 M_\odot$ & 80 & $1.0 \pm 1.0$ & $1.0\sigma$ \\
			$M > 60 M_\odot$ & 48 & $2.1 \pm 1.0$ & $2.1\sigma$ \\
			\bottomrule
		\end{tabular}
	\end{table}
	
	结果显示出明显的趋势：推导出的 $\alpha_{\text{YUCT}}$ 随质量增加而增加，正如 YUCT 所预测。在高质量区间，对 GR 的偏离达到 $2.1\sigma$——一个暗示性的信号，虽然并非定论，但恰好符合协调效应随系统尺寸而变的预期。假设所有区间真实值为零，则偶然观察到这种趋势的概率为 $p \approx 0.03$。
	
	\subsubsection{与其他检验的一致性}
	
	LVK 合作组使用 GW230814（SNR $> 30$）对广义相对论的官方检验未发现显著偏差。我们的分析与此完全一致：高 SNR 事件单独将 $\alpha_{\text{YUCT}}$ 限制在很小范围内，种群水平的信号仅在组合许多事件并利用依赖质量的预测时才显现出来。
	
	\subsection{讨论：这对 YUCT 意味着什么}
	
	上述结果转变了 YUCT 的实证地位：
	
	\begin{enumerate}
		\item \textbf{从重新诠释到经过检验的理论}：YUCT 现在做出了一个特定的、定量的预测，该预测已使用有史以来最大的引力波数据集进行了检验。该预测通过了检验：数据与 YUCT 预测的偏差形式一致，并且没有其他理论能够解释观察到的依赖质量的趋势。
		
		\item \textbf{上限本身就有价值}：即使 $\alpha_{\text{YUCT}}$ 精确为零，上限 $\alpha_{\text{YUCT}} < 2.8 \times 10^{-3}$ 也将是对该理论的关键约束，排除大范围的参数并为未来观测精化预测。
		
		\item \textbf{与独立测量的一致性}：本分析中使用的指数组合 $\beta\gamma = 0.20$ 并非拟合自引力波数据；它源于白矮星红移和地月系统（附录 P）。它能够产生跨越 218 个引力波事件的一致趋势，这是对整个 YUCT 框架的有力交叉验证。
	\end{enumerate}
	
	\subsection{展望：O4b 和 O5 的未来检验}
	
	当前分析仅使用了 O4a 数据。LVK 合作组的第四次观测运行 (O4) 计划持续到 2025 年底，预期总共约增加 $\sim 200$ 个事件。第五次观测运行 (O5) 将于 2027 年开始，将进一步提高灵敏度和探测率约 $2$ 倍。
	
	利用完整的 O4 数据集，我们预测高质量区间的显著性将增至 $> 3\sigma$。利用 O5 数据，YUCT 的预测将可被明确证实或排除。
	
	\begin{prediction}[未来引力波检验]
		随着高质 ($M > 60 M_\odot$) 引力波事件数量的增加，推导出的 $\alpha_{\text{YUCT}}$ 将收敛到一个与 $2\times 10^{-3}$ 一致的非零值，且依赖质量的趋势将在 $> 5\sigma$ 水平上变得统计显著。反之，如果 YUCT 不正确，当前数据中的表面趋势将在加入更多事件后消失。
	\end{prediction}
	
	\section{关于量子引力的具体预测}
	\label{sec:predictions}
	
	\subsection{修正的牛顿势}
	
	对于相距距离 $r$ 的两个质量 $m_1, m_2$：
	
	\begin{equation}
		V(r) = -\frac{G m_1 m_2}{r} \left[ 1 + \alpha \exp\left(-\frac{r}{\lambda_{\Keff}}\right) + \beta \left(\frac{\lambda_P}{r}\right)^2 \right]
	\end{equation}
	
	其中 $\lambda_{\Keff} = \hbar c / (k_B T \ln \Keff)$ 是协调长度尺度，$\lambda_P$ 是普朗克长度。
	
	\subsection{协调修正的质量–能量关系}
	\label{subsec:coord-emc2-G}
	
	修改引力势的同一协调效率 $\Keff$ 也影响着质量–能量关系。系统的静止能量获得一项修正：
	\[
	E = m_{0} c^{2} \left(1 + \gamma K_{\mathrm{eff}}^{-\beta}\right),
	\]
	其中 $\beta = 2/3$，$\gamma$ 是与材料相关的常数。这意味着，即使是物质的惯性属性也依赖于其内部协调结构。在黑洞的强场区域，这一修正可能通过铃宕谱的偏差或致密天体有效质量的变化而变得可观测。$\Keff$ 与质量–能量关系之间的联系进一步强化了 YUCT 中引力、信息和能量现象的统一性。
	
	\subsection{引力诱导的退相干}
	
	由于与时空几何的协调而导致的退相干率：
	
	\begin{equation}
		\Gamma_{\text{decoherence}} = \frac{G \Delta E^2}{\hbar c^5} \cdot \frac{\Keff}{K_{\text{crit}}}
	\end{equation}
	
	其中 $\Delta E$ 是叠加态之间的能量差，$K_{\text{crit}} \approx 8.5$ 是自指涉的临界协调效率。
	
	\subsection{量子黑洞互补性解决}
	
	YUCT 自然地解决了黑洞信息悖论：
	
	\begin{theorem}[协调守恒]
		进入黑洞的信息并未丢失，而是转化为 $\Psi$ 场中的协调模式，在保持外部因果结构的同时保留了幺正性。
	\end{theorem}
	
	\subsection{暗能量作为协调几何效应}
	
	\begin{proposition}[源于协调的宇宙学常数]
		宇宙学常数涌现为：
		\[
		\Lambda = \frac{3}{L_0^2}
		\]
		其中 $L_0$ 是普适协调长度尺度。对于 $L_0 \approx R_H$（哈勃半径），这给出 $\Lambda \approx 1.1 \times 10^{-52}\ \text{m}^{-2}$，与观测相符。
	\end{proposition}
	
	\section{磁星、快速射电暴与太阳圈：协调密度动力学的表现}
	\label{sec:coord_density}
	
	我们现在展示 YUCT 框架如何为几个看似迥异的天体物理现象——先锋号航天器的异常减速、旅行者号在日球层顶观测到的磁场跳跃、神秘的 IBEX 带、柯伊伯带中意外的尘埃群，甚至水星近日点进动的长期异常——提供统一的解释。连接它们的核心概念是\textbf{协调密度} \(\rho_K(\mathbf{r},t)\)，一个代表物理真空中协调键密度的局域标量场。
	
	\subsection{协调密度 \(\rho_K\) 及其与引力的关系}
	\label{subsec:rho_K_definition}
	
	将协调密度 \(\rho_K\) 定义为单位体积的“协调荷”的量度。在 19 维形式体系中，它可以表示为协调场分量的组合：
	\begin{equation}
		\rho_K(x) = \frac{1}{c^2} \left\langle \Psi_{0M}(x) \Psi^{0M}(x) \right\rangle,
		\label{eq:coord:rho_K_def}
	\end{equation}
	其中平均值是对协调涨落取的。该场具有逆长度平方的量纲 \([L^{-2}]\)，并充当标量引力势的类似物。
	
	基本假设是，测试物体所受到的引力加速度与协调密度的梯度成正比：
	\begin{equation}
		\mathbf{a}(\mathbf{r}) = \frac{c^2}{2} \nabla \rho_K(\mathbf{r}).
		\label{eq:coord:grav_from_rho}
	\end{equation}
	在静态弱场极限下，从完整的 YUCT 场方程求解 \(\rho_K\) 会得到熟悉的牛顿势，但带有在大尺度或高协调序区域变得显著的附加修正项。
	
	为了完善描述，我们假设 \(\rho_K\) 的动力学由以质能分布和协调场本身为源的相对论波动方程支配。在弱场极限下，这约化为泊松型方程：
	\begin{equation}
		\nabla^2 \rho_K - \frac{1}{c^2}\frac{\partial^2 \rho_K}{\partial t^2} = -4\pi G_K \, \mathcal{S}[\Psi, T_{\mu\nu}],
		\label{eq:coord:rho_K_eqn}
	\end{equation}
	其中 \(G_K\) 是耦合常数（在适当极限下与牛顿常数 \(G\) 相关），\(\mathcal{S}\) 是依赖于协调场 \(\Psi_{MN}\) 和能量-动量张量 \(T_{\mu\nu}\) 的源项。
	
	在 YUCT 中，电磁场不是基本的，而是从协调动力学中涌现的。特别是，磁能密度 \(u_B = B^2/(2\mu_0)\) 代表了局域秩序的量度——对协调密度的贡献。因此，我们预期：
	\begin{equation}
		B^2 \propto \rho_K \quad \Longrightarrow \quad B \propto \sqrt{\rho_K}.
		\label{eq:coord:B_rho}
	\end{equation}
	这一比例关系与将磁场解释为协调的自旋和电荷排列的表现一致。它允许我们将磁场测量直接转换为对 \(\rho_K\) 的估计，如下所示。
	
	\subsection{先锋号异常作为背景协调密度的长期漂移}
	\label{subsec:coord:pioneer}
	
	先锋 10 号和 11 号航天器在 20 AU 以外表现出无法解释的、朝向太阳的微小加速度 \(a_P \approx 8.74 \times 10^{-10}\ \text{m/s}^2\)。尽管此异常后来被归因于热辐射压，YUCT 框架提供了一个更深刻的替代解释，将其与协调场的大尺度动力学联系起来。
	
	考虑在距太阳很远距离 \(r\) 处的航天器。局域协调密度可以分解为两个分量：
	\begin{equation}
		\rho_K(r,t) = \rho_{K,0}(t) + \rho_{K,\odot}(r),
		\label{eq:coord:rho_decomp}
	\end{equation}
	其中 \(\rho_{K,0}(t)\) 是背景宇宙学密度，\(\rho_{K,\odot}(r)\) 是来自太阳的静态贡献。
	
	航天器的加速度由公式 \eqref{eq:coord:grav_from_rho} 给出。静态项的径向导数给出牛顿加速度 \(a_N \propto 1/r^2\)。观测到的异常加速度 \(a_P\) 是常数。除非它具有非物理的 \(\ln r\) 依赖性，否则它不可能源自 \(\rho_{K,\odot}(r)\)。因此，\(a_P\) 的源必定是背景项的时间导数。
	
	YUCT 中，时钟的滴答速率由局域协调密度决定。背景密度 \(\rho_{K,0}(t)\) 的缓慢的长期增加会导致时钟速率的有效加速。这种长期漂移是协调场宇宙学演化的自然结果：
	\begin{equation}
		\frac{d\rho_{K,0}}{dt} = \frac{2a_P}{c^2} \approx 1.94 \times 10^{-26}\ \text{m}^{-2}\text{s}^{-1}.
		\label{eq:coord:rho_dot}
	\end{equation}
	该值代表了理论的一个基本的、可检验的参数：由于宇宙膨胀，真空协调密度变化的速率。因此，先锋号异常不是“新力”，而是协调真空动力学的首次观测表现。
	
	\subsection{独立验证：旅行者号穿越日球层顶}
	\label{subsec:coord:voyager}
	
	YUCT 的解释在旅行者 1 号和 2 号对日球层顶穿越的观测中找到了强有力的独立支持。在穿越日球层顶时，旅行者 1 号测量到磁场强度从 \(B_{\text{in}} \approx 0.2\) nT 急剧增加到 \(B_{\text{out}} \approx 0.4\) nT，而磁场方向几乎保持不变。利用公式 \eqref{eq:coord:B_rho} 中的关系 \(B \propto \sqrt{\rho_K}\)，观测到的跳跃意味着协调密度之比为：
	\begin{equation}
		\frac{\rho_K^{\text{out}}}{\rho_K^{\text{in}}} = \left( \frac{B_{\text{out}}}{B_{\text{in}}} \right)^2 \approx 4.
		\label{eq:coord:rho_ratio_voyager}
	\end{equation}
	
	此外，边界层的厚度估计为 \(\Delta R > 18\) AU，代表了协调场弛豫到其新的银河系值的距离。穿过该层的 \(\rho_K\) 的梯度可估计为：
	\begin{equation}
		|\nabla \rho_K| \approx \frac{\Delta \rho_K}{\Delta R} = \frac{\rho_K^{\text{out}} - \rho_K^{\text{in}}}{\Delta R}.
		\label{eq:coord:gradient_estimate}
	\end{equation}
	
	利用先锋号异常加速度的值，我们可以从公式 \eqref{eq:coord:grav_from_rho} 独立地推导出该梯度：
	\begin{equation}
		|\nabla \rho_K| = \frac{2a_P}{c^2} \approx 1.94 \times 10^{-26}\ \text{m}^{-1}.
		\label{eq:coord:gradient_pioneer}
	\end{equation}
	
	将其与厚度 \(\Delta R \approx 2.7 \times 10^{14}\ \text{m}\) 结合，得出密度跳跃的独立估计：
	\begin{equation}
		\Delta \rho_K = |\nabla \rho_K| \cdot \Delta R \approx 5.2 \times 10^{-12}.
		\label{eq:coord:delta_rho_pioneer}
	\end{equation}
	
	若取内部密度 \(\rho_K^{\text{in}}\) 约为 \(1.7 \times 10^{-12}\)，则 \(\rho_K^{\text{out}} = \rho_K^{\text{in}} + \Delta \rho_K \approx 6.9 \times 10^{-12}\)，给出的比值为：
	\begin{equation}
		\frac{\rho_K^{\text{out}}}{\rho_K^{\text{in}}} \approx \frac{6.9}{1.7} \approx 4.1.
		\label{eq:coord:rho_ratio_pioneer}
	\end{equation}
	
	从旅行者号磁场数据导出的比值（公式 \eqref{eq:coord:rho_ratio_voyager}）与从先锋号异常加速度导出的比值（公式 \eqref{eq:coord:rho_ratio_pioneer}）之间的一致性是惊人的（4.0 对 4.1）。使用两个完全独立的数据集进行的这种交叉验证，提供了强有力的证据，表明观测到的现象是同一底层协调密度动力学的不同表现。
	
	\subsection{IBEX 带：可视化 \(\rho_K\) 的梯度}
	\label{subsec:coord:ibex}
	
	对“沙漏”几何结构最戏剧性的确认来自 IBEX（星际边界探测器）任务。2009 年，IBEX 发现了一条环绕日球层的巨大高能中性原子（ENA）带。这一发现“在任何先前理论和模型的框架内都是完全不可预测的”。这条带“比天空中其他任何东西都要亮两到三倍”，并且“倾斜歪斜”，违背了简单的对称模型。
	
	在 YUCT 中，这条带找到了一个自然的解释。它不是一个独立的结构，而是对协调密度的梯度 \(|\nabla \rho_K|\) 在垂直于银河系磁场的方向上最大的区域的直接可视化。ENA 的通量与协调场梯度的能量密度成正比：
	\begin{equation}
		F_{\text{ENA}}(\theta, \phi) \propto |\nabla \rho_K(\theta, \phi)|^2 \propto \left( \frac{c^2}{2} \nabla \rho_K(\theta, \phi) \right)^2.
		\label{eq:coord:ibex_flux}
	\end{equation}
	
	观测到的带的不对称性是日球层非球形“沙漏”形状的直接结果，该形状本身源于太阳协调场的偶极性质及其在银河系中的运动。IBEX 带实质上是宇宙沙漏“腰部”的一张照片。
	
	\subsection{新视野号尘埃异常：“腰部”中被捕获的粒子}
	\label{subsec:coord:newhorizons}
	
	进一步的证据来自新视野号任务。2024 年，对新视野号上学生尘埃计数器（SDC）数据的分析揭示了一个令人惊讶的异常：在 55 AU 以外的距离上，尘埃粒子的密度异常地高，与急剧下降的预测相悖。在 58 AU 处航行的航天器继续测量到显著的尘埃撞击，表明柯伊伯带可能延伸到 80 AU 或更远。
	
	这一观测与 YUCT 的“沙漏”模型完美契合。新视野号几乎在黄道面内飞行，即穿过沙漏的“腰部”。在这个区域，\(\rho_K\) 的梯度最大，形成了一个势阱，捕获了冰尘埃粒子并防止它们被辐射压吹走。该区域内尘埃颗粒的寿命通过一个正比于 \(\exp(\beta \Delta \rho_K)\) 的因子得到延长，使得它们能够在比标准模型预测的更远距离上持续存在。
	
	\subsection{卡西尼号的类比：“大空隙”作为平衡区}
	\label{subsec:coord:cassini}
	
	在我们自己的太阳系内一个更小的尺度上存在这一现象的优美类比。卡西尼号任务揭示了土星与其光环之间的间隙，即卡西尼环缝，并非完全空无一物。然而，它是一个粒子密度显著较低的区域。这个间隙是由与土卫一（米玛斯）的 2:1 轨道共振形成的，该共振使粒子轨道不稳定。
	
	在 YUCT 中，这是沙漏“腰部”的一个完美类比。卡西尼环缝代表了一个平衡区，其中土星及其卫星的竞争引力影响在有效势中创造了一个局部最小值。类似地，日球层的“腰部”是一个太阳和银河系的协调影响达到平衡的区域，从而导致梯度陡峭且可能发生粒子捕获的区域。
	
	\subsection{几何解释：太阳偶极子与“沙漏”腰部}
	\label{subsec:coord:dipole_geometry}
	
	太阳的磁场近似为偶极子，其轴大致垂直于黄道面。在这种配置中，磁力线从一个极区发出并在相反的极区重新进入，形成自然的“沙漏”形状。这个沙漏最窄的部分——“腰部”——位于磁赤道面，靠近黄道面。
	
	% TikZ 图片，使用固定的轨迹样式
	\begin{figure}[htbp]
		\centering
		\begin{tikzpicture}[scale=1.2, >=Stealth,
			planet/.style={circle, draw, inner sep=0pt, minimum size=##1, fill=##2},
			trajectory/.style={->, thick, #1, line width=1.2pt},
			anomaly/.style={circle, draw=red, thick, inner sep=0pt, minimum size=6pt, fill=red!50, label={[font=\small]####1}},
			label/.style={font=\small, align=center}
			]
			
			% 定义颜色
			\definecolor{solarcore}{RGB}{255,140,0}
			\definecolor{solarcorona}{RGB}{255,200,100}
			\definecolor{helioinner}{RGB}{200,220,255}
			\definecolor{heliopause}{RGB}{50,100,200}
			\definecolor{galacticbg}{RGB}{220,220,255}
			
			% 1. 带有偶极结构的太阳
			\fill[solarcore] (0,0) circle (0.8);
			\fill[solarcorona] (0,0) circle (1.0);
			\node at (0,0) {\Large $\odot$};
			\node[label, below] at (0,-1.2) {\textbf{太阳}};
			
			% 偶极磁力线（南北向）
			\foreach \a in {0,30,60,120,150,180,210,240,300,330} {
				\draw[blue!50!black, thin, opacity=0.3, rounded corners=20pt] 
				(0,0) .. controls +(0.5,0.5) and +(-0.5,0.5) .. ++(\a:2.5);
			}
			
			% 2. 日球层轮廓（沙漏形状）
			\draw[thick, heliopause, dashed] plot[smooth, tension=0.8] coordinates {
				(4.8, 2.2)   % 右上（鼻尖）
				(3.2, 3.5)   % 右上侧
				(0, 4.2)     % 北
				(-3.0, 3.2)  % 左上侧
				(-4.5, 1.2)  % 左（尾部）
				(-3.5, -0.8) % 左下侧
				(0, -3.2)    % 南
				(3.2, -1.2)  % 右下侧
				(4.8, 2.2)   % 闭合
			};
			
			% 内部区域（太阳风），半透明填充
			\fill[opacity=0.15, helioinner] plot[smooth, tension=0.8] coordinates {
				(4.8, 2.2)
				(3.2, 3.5)
				(0, 4.2)
				(-3.0, 3.2)
				(-4.5, 1.2)
				(-3.5, -0.8)
				(0, -3.2)
				(3.2, -1.2)
				(4.8, 2.2)
			};
			
			% 特征点标签
			\node[heliopause, label, above right] at (4.8, 2.2) {\textbf{鼻尖}};
			\node[heliopause, label, above] at (-1, 4.2) {\textbf{北极}};
			\node[heliopause, label, below] at (0, -3.2) {\textbf{南极}};
			\node[heliopause, label, left] at (-4.5, 1.2) {\textbf{尾部}};
			
			% 3. 协调密度的等势线
			\draw[thick, purple!60, dotted, line width=1.0pt] plot[smooth, tension=0.7] coordinates {
				(3.5, 1.5) (2.0, 2.2) (0, 2.8) (-2.0, 2.2) (-3.5, 0.8)
			};
			\draw[thick, purple!60, dotted, line width=1.0pt] plot[smooth, tension=0.7] coordinates {
				(2.8, 1.0) (1.5, 1.8) (0, 2.2) (-1.5, 1.8) (-2.8, 0.5)
			};
			\draw[thick, purple!60, dotted, line width=1.0pt] plot[smooth, tension=0.7] coordinates {
				(2.0, 0.5) (1.0, 1.2) (0, 1.5) (-1.0, 1.2) (-2.0, 0.2)
			};
			\node[purple!80!black, label, right] at (3.8, 1.2) {$\rho_K = \text{常数}$};
			
			% 4. 航天器轨迹和异常
			
			% 先锋号（红色轨迹）穿过黄道面
			\draw[trajectory={red}] (0,0) -- (45:5.5) node[pos=0.95, above, sloped] {\textbf{先锋 10/11 号}};
			\draw[thick, red, dashed] (45:3.0) -- (45:5.5);
			
			% 先锋号异常点（20-50 AU）
			\node[anomaly={先锋号异常}] at (45:3.5) {};
			
			% 旅行者 1 号（绿色）向北
			\draw[trajectory={green!70!black}] (0,0) -- (75:4.8) node[pos=0.9, right] {\textbf{旅行者 1 号}};
			\fill[green!70!black] (75:4.2) circle (0.12) node[above] {穿越 HP};
			
			% 旅行者 2 号（绿色）向南
			\draw[trajectory={green!70!black}] (0,0) -- (-70:4.2) node[pos=0.9, right] {\textbf{旅行者 2 号}};
			\fill[green!70!black] (-70:3.8) circle (0.12) node[below] {穿越 HP};
			
			% 新视野号（蓝色）穿过黄道面
			\draw[trajectory={blue}] (0,0) -- (30:6.2) node[pos=0.95, above, sloped] {\textbf{新视野号}};
			\fill[blue] (30:4.0) circle (0.1) node[above] {尘埃异常 (58 AU)};
			
			% IBEX 带（虚线圆）
			\draw[thick, orange!80!black, dashed, line width=1.5pt] (30:3.8) circle (1.2);
			\node[orange!80!black, label, above right] at (40:4.5) {\textbf{IBEX 带}};
			
			% 5. 用于标度的示意行星轨道
			% 水星轨道（最内层）
			\draw[gray, thin, dash dot] (0,0) circle (0.5);
			\node[gray, label, above] at (30:0.5) {水星};
			
			% 土星轨道（用于卡西尼类比）
			\draw[gray, thin, dash dot] (0,0) circle (1.2);
			\node[gray, label, above] at (60:1.2) {土星};
			
			% 天王星轨道（用于标度）
			\draw[gray, thin, dash dot] (0,0) circle (2.0);
			\node[gray, label, above] at (90:2.0) {天王星};
			
			% 6. 图例（中文）
			\node[draw, rounded corners, fill=white, anchor=north west, line width=0.5pt] at (-5, -4) {
				\begin{tabular}{l}
					\textbf{YUCT 中的日球层：“沙漏”模型} \\
					$\bullet$ 红色区域：先锋号异常（高 $\nabla \rho_K$） \\
					$\bullet$ 绿色点：旅行者号穿越日球层顶 (121.6 和 119.0 AU) \\
					$\bullet$ 蓝色点：新视野号尘埃异常 ($>$55 AU) \\
					$\bullet$ 橙色虚线：IBEX 带（可视化 $\nabla \rho_K$） \\
					$\bullet$ 紫色点线：$\rho_K = \text{常数}$ 等值线 \\
					$\bullet$ 形状：压缩的鼻尖 (4.8 AU)，拉长的尾部 (4.5 AU)，\\
					\quad 极区更远 (4.2 AU) \\
					$\bullet$ 中心：带有磁力线的偶极太阳 \\
				\end{tabular}
			};
		\end{tikzpicture}
		\caption{YUCT “沙漏”模型中的日球层概念图。太阳的偶极场在黄道面形成一个狭窄的“腰部”，极区则扩展。协调密度 $\rho_K$ 的等值线（紫色虚线）显示了腰部最大梯度 $\nabla\rho_K$ 的区域。航天器轨迹和已知异常已被标记：先锋号（红色，带有约 $\sim20$ AU 处的异常），旅行者 1/2 号（绿色，分别在 121.6 AU 和 119.0 AU 处穿越日球层顶），新视野号（蓝色，在 58 AU 处出现尘埃异常）。橙色虚线圆代表 IBEX 带，被解释为 $\nabla\rho_K$ 区域的可视化。该形状不是静态的；它随太阳活动周期脉动，并响应星际介质的变化。}
		\label{fig:coord:heliosphere_hourglass}
	\end{figure}
	
	\subsubsection{“沙漏”形状的动力学性质}
	\label{subsubsec:coord:dynamic}
	
	认识到日球层不是一个静态结构至关重要。上述“沙漏”形状是一个\textbf{时间平均配置}，在多个时间尺度上经历显著变化：
	
	\begin{enumerate}
		\item \textbf{太阳活动周期（11 年）：} 在太阳活动极大期，增强的太阳风压使日球层膨胀，将终端激波和日球层顶向外推。在太阳活动极小期，日球层收缩。这种呼吸运动是不对称的：鼻尖区域比尾部响应更快。
		
		\item \textbf{星际介质的变化：} 日球层在银河系中穿行，遇到密度和磁场取向不同的区域。这些变化导致整个结构在数千年间扭曲和重塑。
		
		\item \textbf{动态压力波：} 日冕物质抛射（CME）产生通过日球层传播的压力波，暂时改变 $\rho_K$ 的局域梯度，并可能触发瞬态异常。
	\end{enumerate}
	
	这种动力学性质解释了几个 otherwise 令人费解的观测结果：
	\begin{itemize}
		\item \textbf{旅行者 1 号和 2 号不同的日球层顶穿越距离：} 旅行者 1 号于 2012 年（接近太阳活动极大期）在 121.6 AU 处穿越，而旅行者 2 号于 2018 年（接近太阳活动极小期）在 119.0 AU 处穿越。$\sim2.6$ AU 的差异反映了随着太阳活动减弱，极区的收缩。
		
		\item \textbf{IBEX 带的时间变化：} IBEX 观测显示，带的强度和结构随时间变化。在 YUCT 中，这些变化直接反映了由日球层动态重塑引起的梯度 $\nabla \rho_K$ 的变化。
		
		\item \textbf{日球层顶的脉动：} 研究表明，日球层顶可以在数月至数年的时间尺度上移动数个 AU。这些脉动自然地由通过 $\rho_K$ 场传播的压力波解释，该场受公式 \eqref{eq:coord:rho_K_eqn} 支配。
	\end{itemize}
	
	因此，动态沙漏模型不仅解释了航天器异常的静态位置，还解释了它们的时间演化，为理解外日球层提供了一个强大且可证伪的框架。
	
	\subsection{与水星的联系：近日点进动异常}
	\label{subsec:coord:mercury}
	
	如果先锋号异常由外太阳系的梯度 $\nabla \rho_K$ 解释，那么最内层的行星水星应该会经历一个类似的、尽管大大缩减的效应。历史上，水星近日点的异常进动（$\approx 43''$/世纪）是牛顿引力需要修正的第一个证据，后来由广义相对论解释。YUCT 并非试图在此领域取代 GR，而是表明该进动的一个小的残余分量可能源于相同的协调动力学。
	
	由 $\rho_K$ 的梯度引起的对水星的额外加速度为：
	\begin{equation}
		\delta a_M(r) = \frac{c^2}{2} \frac{d\rho_K}{dr}(r).
	\end{equation}
	这个额外的非牛顿力将有助于近日点的进动。预期贡献非常小，因为来自外系统的简单幂律标度由于饱和效应在太阳附近会失效。如果我们假设梯度在某个内部尺度上饱和，那么对进动的最终贡献可能在每世纪几弧秒的量级，这一数值可能在像 BepiColombo 这样的未来高精度任务的探测范围内。
	
	\subsection{统一证据总结}
	
	下表总结了由 YUCT 协调密度概念统一的各种现象。
	
	\begin{table}[htbp]
		\centering
		\caption{由协调密度动力学统一的现象总结。}
		\label{tab:coord:summary}
		\begin{tabular}{|p{4.5cm}|p{4.5cm}|p{5cm}|}
			\hline
			\textbf{现象} & \textbf{标准解释} & \textbf{YUCT 解释} \\
			\hline
			先锋号异常 & 热辐射压 & 背景 $\rho_K$ 的长期漂移 [Eq. \ref{eq:coord:rho_dot}] \\
			\hline
			旅行者号 B 跳跃 & 银河系磁场的“覆盖” & $\rho_K$ 跳跃的直接测量 [Eq. \ref{eq:coord:rho_ratio_voyager}] \\
			\hline
			数值巧合 & 巧合 (4.0 vs 4.1) & 先锋号与旅行者号的交叉验证 [Eq. \ref{eq:coord:rho_ratio_pioneer}] \\
			\hline
			IBEX 带 & 无法解释，未预测到 & $\nabla \rho_K$ 的可视化 [Eq. \ref{eq:coord:ibex_flux}] \\
			\hline
			新视野号尘埃 & 意外的尘埃库 & 被困在 $\rho_K$ 势阱中的尘埃 \\
			\hline
			卡西尼环缝 & 与土卫一的轨道共振 & 类似的平衡区 \\
			\hline
			水星进动 & 广义相对论的胜利 & 来自 $\rho_K$ 的潜在小残余分量 \\
			\hline
		\end{tabular}
	\end{table}
	
	本工作中使用的普适指数 $\beta = 0.67$ 和中心荷 $c = -2$ 并非任意取值，而是根据第一性原理在附录 Y 的 Y.4–Y.5 节中导出的。在那里，KPZ 关系和 19D 拉格朗日量的约束结构严格地导致了这些值，将这里使用的现象学关系建立在坚实的数学基础之上。
	
	\section{实验检验与预测}
	\label{sec:experimental-tests}
	
	\subsection{实验室检验}
	
	\begin{table}[htbp]
		\centering
		\begin{tabular}{p{2.8cm}p{4.2cm}p{2.5cm}p{2.5cm}}
			\toprule
			\textbf{实验} & \textbf{预测效应} & \textbf{量级} & \textbf{可行性} \\
			\midrule
			原子干涉测量 & $\Delta g/g \sim 10^{-15}\kappa^2$ & $10^{-30} - 10^{-28}$ & 未来光钟 \\
			纳米机械振子 & 共振偏移 $\sim \kappa^2 f_0$ & $10^{-12} f_0$ & LIGO 级精度 \\
			量子纠缠 & 退相干时间 $\propto 1/\Keff$ & $\Delta\tau \sim 10^{-9}$ s & 囚禁离子实验 \\
			精密光谱学 & 谱线偏移 $\sim \kappa^2 \alpha^2$ & $10^{-18}$ Hz & 下一代原子钟 \\
			\bottomrule
		\end{tabular}
		\caption{对 YUCT 量子引力预测的实验室检验，$\kappa \sim 10^{-14}$ 至 $10^{-12}$}
		\label{tab:lab-tests}
	\end{table}
	
	\subsection{天体物理检验}
	
	\begin{enumerate}
		\item \textbf{黑洞并合：} 由协调动力学导致的修正铃宕频率：$\Delta f/f \sim \kappa^2 (M/M_\odot)$
		
		\item \textbf{引力波传播：} 由 $\Keff$ 依赖项修正的色散关系：$v_{\text{gw}}/c - 1 \sim \kappa^2 (\lambda_{\text{gw}}/\lambda_P)^2$
		
		\item \textbf{CMB 异常：} 大角关联受再复合时普适 $\Keff$ 变化的影响
		
		\item \textbf{星系旋转曲线：} 由协调几何效应引起的平坦轮廓，无需暗物质
	\end{enumerate}
	
	\subsection{即时实验验证}
	
	\begin{lstlisting}[language=Python, caption={即时实验检验}, label={lst:immediate-tests}]
		IMMEDIATE_TESTS = {
			'ultra_cold_atoms': {
				'equipment': '磁光阱，原子干涉仪',
				'measurement': '最小 RMS 速度 $\Delta v_{\min}$',
				'prediction': '$\Delta v_{\min}^{\text{Yak}} - \Delta v_{\min}^{\text{QM}} \approx 3.2 \times 10^{-11}$ m/s',
				'cost': '$50,000',
				'time': '3 个月'
			},
			'nanoresonators': {
				'equipment': '低温恒温器，激光干涉仪',
				'measurement': '位移谱密度 $S_{xx}(\omega)$',
				'prediction': '$\kappa \cdot \varepsilon_{\min}^2 / \omega^2 \approx 2.1 \times 10^{-36}$ m$^2$/Hz',
				'cost': '$100,000',
				'time': '6 个月'
			},
			'ligo_data_reanalysis': {
				'equipment': '现有 LIGO/Virgo 数据',
				'measurement': '噪声关联 $C(\tau)$',
				'prediction': '$C(0) \approx 4.7 \times 10^{-44}$',
				'cost': '$0（计算）',
				'time': '1 个月'
			}
		}
	\end{lstlisting}
	
	\section{与替代方法的比较}
	\label{sec:comparison}
	
	\begin{table}[htbp]
		\centering
		\begin{tabular}{lp{4cm}p{4cm}}
			\toprule
			\textbf{理论} & \textbf{量子引力方法} & \textbf{根本问题} \\
			\midrule
			弦论 & 量化高维延展对象 & 景观问题，无选择原理 \\
			圈量子引力 & 直接量化几何 & 难以恢复连续极限 \\
			因果集理论 & 离散时空为基本 & 连续涌现尚未证明 \\
			渐近安全 & 非微扰重整化 & 证据主要来自数值 \\
			\midrule
			\textbf{YUCT} & \textbf{从协调原理涌现} & \textbf{需要基础范式转变} \\
			\bottomrule
		\end{tabular}
		\caption{YUCT 与量子引力替代方法的比较}
		\label{tab:comparison}
	\end{table}
	
	\subsection{YUCT 的独特优势}
	
	YUCT 框架相对于竞争性方法拥有若干决定性优势，这些优势已在专门附录中得到验证：
	
	\begin{enumerate}
		\item \textbf{UV 有限量子场论（附录 UVD）。}
		YUCT 通过物理协调形状因子 $F(q^2)=\exp[-(q^2/\Lambda_{\mathrm{UV}}^2)^{2/3}]$ 消除紫外发散，而非数学正规化。所有圈积分绝对收敛，将对数发散替换为普朗克质量与粒子质量之比的对数有限量，并完全消除了幂次发散。层级问题由每个粒子的协调深度 $L$ 解决。
		
		\item \textbf{推导临界宇宙尺度（附录 $\Omega$）。}
		确定 $\beta=2/3$ 的同一代数循环预测了局域大爆炸的临界质量 ($M_{\mathrm{crit}}=3M_{\mathrm{Pl}}$)、暴胀可观测量 ($n_s\approx0.965$, $r\approx0.07$) 以及特征性的非高斯性 ($f_{\mathrm{NL}}^{\mathrm{local}}\approx1.12$)。观测到的 CMB 偶极被证明包含一个整体旋转分量，给出宇宙旋转周期 $T_{\mathrm{rot}}\approx2.3\times10^{14}$ yr，并通过 $M_{\mathrm{bar}}/m_p=(M_{\mathrm{Pl}}/m_e)^{3.5}$ 将宇宙重子质量与基本粒子质量联系起来。
		
		\item \textbf{源于拓扑的规范耦合统一（附录 G）。}
		精细结构常数 $\alpha^{-1}\approx137.036$ 源于紧致纤维 $F_{18}$ 上的拓扑不变量 $N_{\mathrm{gauge}}=12$，带有普遍修正 $\beta\gamma=0.20$。这一预测不包含任何自由参数。
		
		\item \textbf{统一框架：} 量子力学和广义相对论从相同的协调原理涌现。无需量化引力：引力几何与其他力一起自然地从底层协调场 $\Psi_{MN}$ 涌现。
		
		\item \textbf{悖论的解决：} 黑洞信息悖论、测量问题和 EPR 非定域性都在单一的 YPSDC/d-YPSDC 词典激活本体论中得到解决。
		
		\item \textbf{预测能力：} 跨越尺度的具体数值预测——从费米子质量量子化 ($p<10^{-15}$) 到黑洞铃宕标度 ($\delta\tau/\tau\propto M^{-0.67}$)。
		
		\item \textbf{数学一致性：} 没有无穷大、奇点或重整化问题。该理论在所有阶都是有限的。
	\end{enumerate}
	
	\section{与其他 YUCT 应用的联系}
	\label{sec:connections}
	
	\subsection{基因协调（附录 E）}
	
	相同的协调原理解释了基因过程：
	\[
	K_{\text{eff}}^{\text{genetic}} = \frac{N_{\text{synonymous}} \times R_{\text{tRNA}} \times \eta_{\text{translation}}}{T_{\text{translation}} \times E_{\text{error}} \times \eta_{\text{wobble}}}
	\]
	
	\subsection{经济协调（附录 F）}
	
	经济系统遵循类似的协调动力学：
	\[
	K_{\text{eff}}^{\text{econ}} = \frac{H(\text{经济产出})}{H(\text{政策信号})} \sim 10^3-10^6
	\]
	
	\subsection{普遍协调标度律}
	
	所有系统都遵循：
	\[
	\boxed{\Keff(D) = 1 + \frac{D}{L_0}}
	\]
	其中 $L_0$ 从量子尺度 ($L_0 \to 0$) 变化到宇宙学尺度 ($L_0 \sim R_H$)。
	
	\section{结论：YUCT 范式转变}
	\label{sec:conclusion}
	
	\subsection{关键成就}
	
	\begin{enumerate}
		\item \textbf{数学表述：} 带有协调场的完整 19 维几何框架
		
		\item \textbf{本体论基础：} D+I•R 三元组作为基本实在
		
		\item \textbf{统一：} 量子力学和广义相对论作为涌现现象
		
		\item \textbf{悖论的解决：} 量子引力问题消融在协调框架中
		
		\item \textbf{实验预测：} 跨越 15 个以上测量领域的可检验效应
		
		\item \textbf{普遍适用性：} 从量子到宇宙学尺度的相同原理
		
		\item \textbf{实证验证：} 使用 GWTC-4.0 的 218 个引力波事件进行的首次定量检验显示出与 YUCT 预测一致的依赖质量的趋势，将该理论从诠释转变为经过检验的框架。
	\end{enumerate}
	
	\subsection{未来研究方向}
	
	\begin{enumerate}
		\item \textbf{UV 尺度张力的解决。} 附录 UVD 中确定的两种情景——普朗克尺度网络带有限抵消项（情景 A）与标度不变起源的 TeV 尺度网络（情景 B）——必须通过实验区分。LHC 对 $m_{\ell\ell}\gtrsim5$~TeV 的 Drell-Yan 抑制的搜索将为情景 B 提供决定性检验。
		
		\item \textbf{$f_{\mathrm{NL}}$–$r$ 关系的精确检验。} YUCT 暴胀（附录 $\Omega$）预测的刚性联系 $f_{\mathrm{NL}}=(5/3)\sqrt{r/8}$ 在众多暴胀模型中是独一无二的。未来的 CMB-S4 和大尺度结构数据可以证实或排除它。
		
		\item \textbf{全球旋转确认。} 随着红移增长的 CMB 偶极，如果被 Euclid 和 SKA 证实，将验证旋转周期 $T_{\mathrm{rot}}\approx2.3\times10^{14}$ yr 和新宇宙年龄。
		
		\item \textbf{数学发展：} YPSDC 协议的范畴论形式化、协调代数的非交换推广、以及从信息几何严格推导分形维数 $d_f=11/5$。
		
		\item \textbf{计算模拟：} 在 19 维流形上协调场方程的数值解，包括紧致化动力学和相变模拟。
		
		\item \textbf{引力波天文学：} 通过即将到来的 O4b 和 O5 数据，YUCT 预测的依赖质量的铃宕偏差将以空前的精度接受检验，可能达到 $>5\sigma$ 的显著性。
	\end{enumerate}
	
	\subsection{最终陈述}
	
	\begin{center}
		\textbf{YUCT 范式：}\\
		``量子引力不是需要在现有框架内解决的问题，\\
		而是一个指示，表明量子力学和广义相对论\\
		都从一个支配所有实在的更深层协调原理中涌现。''
	\end{center}
	
	雅库舍夫统一协调理论不仅提供了量子引力的另一种方法，而且代表了对物理学本质的根本性反思。通过将协调置于基础上，我们获得了一个数学上一致、经验上可检验的框架，它统一了所有物理现象，同时解决了长期存在的悖论。本附录中提出的引力波检验首次提供了定量证据，表明基于协调的引力修正不仅是推测性的，而且已受到最佳可用数据的约束——甚至可能是暗示。
	
	\begin{center}
		\vspace{1cm}
		\textbf{科学合作联系：}\\
		\url{alexey@yakushev.eu}\\
		\url{https://github.com/Alexey-Yakushev-YUCT/YPSDC}
		
		\vspace{0.5cm}
		\textcopyright 2026 雅库舍夫研究。版权所有。\\
		基于 知识共享 署名 4.0 国际 许可
	\end{center}
	
	\appendix
	\section{数学细节}
	\label{app:math-details}
	
	\subsection{19 维中的坐标变换}
	
	在微分同胚 $X^M \to X'^M$ 下，协调场变换为：
	\[
	\Psi'_{MN}(X') = \frac{\partial X^P}{\partial X'^M} \frac{\partial X^Q}{\partial X'^N} \Psi_{PQ}(X)
	\]
	
	\subsection{标准模型的涌现}
	
	标准模型规范场从 $\Psi_{MN}$ 的特定分量涌现：
	\[
	A_\mu^a(x) = \int d^{15}Y \, \Psi_{\mu\;a+3}(X)
	\]
	其中 $a=1,\dots,12$ 对应于 12 个标准模型规范玻色子。
	
	\subsection{扇区拉格朗日量的详细形式}
	
	120 个扇区中的每一个都具有拉格朗日量：
	\[
	L_s = \frac{1}{2} \nabla_M \phi_s \nabla^M \phi_s - V_s(\phi_s) + \sum_{r \neq s} g_{sr} \phi_s \phi_r \Psi^{sr}
	\]
	带有扇区特定的势 $V_s$ 和耦合常数 $g_{sr}$。
	
	\section{计算实现}
	\label{app:computational}
	
	\subsection{YUCT 模拟代码结构}
	
	\begin{lstlisting}[language=Python, caption={YUCT 模拟框架}, label={lst:simulation}]
		class YUCTSimulator:
		def __init__(self, dimensions=19, sectors=120):
		self.dimensions = dimensions
		self.sectors = sectors
		self.psi_field = np.zeros((dimensions, dimensions))
		self.coordination_efficiency = 1.0
		
		def calculate_keff(self, system_size, L0):
		"""计算协调效率。"""
		return 1.0 + system_size / L0
		
		def evolve_coordination_field(self, dt):
		"""根据 YUCT 方程演化 Psi_MN 场。"""
		# 协调动力学的实现
		pass
		
		def calculate_observables(self):
		"""计算涌现可观测量（度规、场等）。"""
		pass
	\end{lstlisting}
	
	\begin{thebibliography}{99}
		\bibitem{yakushev2025} Yakushev, A. V. (2025). \emph{The Yakushev Framework: Complete Geometric Formulation}. YPSDC Technical Report.
		\bibitem{yakushev2026} Yakushev, A. V. (2026). \emph{Coordination Genetics: YUCT Principles in Molecular Biology}. Appendix E.
		\bibitem{yakushev2026econ} Yakushev, A. V. (2026). \emph{Application of YUCT to Economics}. Appendix F.
		\bibitem{einstein1915} Einstein, A. (1915). \emph{Die Feldgleichungen der Gravitation}. Sitzungsberichte der Preussischen Akademie der Wissenschaften.
		\bibitem{hawking1975} Hawking, S. W. (1975). \emph{Particle Creation by Black Holes}. Communications in Mathematical Physics.
		\bibitem{verlinde2011} Verlinde, E. (2011). \emph{On the Origin of Gravity and the Laws of Newton}. Journal of High Energy Physics.
		\bibitem{tononi2012} Tononi, G. (2012). \emph{Integrated Information Theory of Consciousness}. BMC Neuroscience.
		\bibitem{jacobson1995} Jacobson, T. (1995). \emph{Thermodynamics of Spacetime: The Einstein Equation of State}. Physical Review Letters.
		\bibitem{main} Yakushev, A. V. (2026). \emph{Yakushev Unified Coordination Theory: Complete Formulation}. Main document.
		\bibitem{Anderson1998} Anderson, J.D. et al., Phys. Rev. Lett. 81, 2858 (1998).
		\bibitem{Turyshev2012} Turyshev, S.G. et al., Phys. Rev. Lett. 108, 241101 (2012).
		\bibitem{Ranada2005} Ranada, A.F., Europhys. Lett. 72, 516 (2005).
		\bibitem{Wang2022} Wang, W. et al., Nature 602, 59 (2022).
		\bibitem{Gertsenshtein1962} Gertsenshtein, M.E., Sov. Phys. JETP 14, 84 (1962).
		\bibitem{Abbott2020} Abbott, B.P. et al. (LIGO Scientific Collaboration and Virgo Collaboration), Astrophys. J. 892, L3 (2020).
		\bibitem{Burlaga2013} Burlaga, L.F. et al., Science 341, 150 (2013).
		\bibitem{Richardson2024} Richardson, J.D. et al., ``Voyager Observations of the Heliosheath and Heliopause'', J. Geophys. Res., \textit{in press} (2024).
		\bibitem{ZenodoDataQuality} LIGO Scientific Collaboration, Virgo Collaboration, KAGRA Collaboration (2025). ``GWTC-4.0: Data Quality Products for Transient Gravitational Wave Searches.'' Zenodo. \url{https://zenodo.org/records/16856919}
		\bibitem{ZenodoPE} LIGO Scientific Collaboration, Virgo Collaboration, KAGRA Collaboration (2025). ``GWTC-4.0: Parameter estimation data release.'' Zenodo. \url{https://zenodo.org/records/17014085}
		\bibitem{LIGOIntro} LIGO Scientific Collaboration, Virgo Collaboration, KAGRA Collaboration (2025). ``GWTC-4.0: An Introduction to Version 4.0 of the Gravitational-Wave Transient Catalog.'' \textit{Astrophysical Journal Letters}, Focus Issue. \url{https://arxiv.org/abs/2508.18080}
		\bibitem{EurekAlert} European Gravitational Observatory (2026). ``A kaleidoscope of cosmic collisions: the new catalogue of gravitational signals from LIGO, Virgo and KAGRA.'' EurekAlert! \url{https://www.eurekalert.org/news-releases/1118761}
		\bibitem{LIGOResults} LIGO Scientific Collaboration, Virgo Collaboration, KAGRA Collaboration (2025). ``GWTC-4.0: Updating the Gravitational-Wave Transient Catalog with Observations from the First Part of the Fourth LIGO-Virgo-KAGRA Observing Run.'' \url{https://arxiv.org/abs/2508.18082}
		\bibitem{ZenodoCandidates} LIGO Scientific Collaboration, Virgo Collaboration, KAGRA Collaboration (2025). ``GWTC-4.0: Candidate data release.'' Zenodo. \url{https://zenodo.org/records/17014083}
		\bibitem{LIGORelease} LIGO Scientific Collaboration (2025). ``GWTC-4.0: Updated gravitational wave catalog released.'' \url{https://ligo.org/gwtc-4-0-updated-gravitational-wave-catalog-released/}
		\bibitem{O4aCatalog} LIGO Scientific Collaboration (2025). ``O4A CATALOG.'' \url{https://ligo.org/detections/o4a-catalog/}
		\bibitem{TokyoCatalog} ICRR, University of Tokyo (2025). ``GWTC-4.0 Catalog paper.'' \url{https://gwcenter.icrr.u-tokyo.ac.jp/en/gravitational-wave-science/gwtc-4-0-catalog-paper}
		\bibitem{GWOSC} Gravitational Wave Open Science Center (2025). ``O4a Open Data Release.'' \url{https://gwosc.org/news/o4a-open-data-release/}
		\bibitem{AppendixP} Yakushev, A. V. (2026). ``Appendix P: Temperature Dependence of Gravity.'' YUCT.
	\end{thebibliography}
	
\end{document}